\documentclass[11pt]{article}
\usepackage{amsmath,amssymb,amsthm,hyperref}
\newtheorem{theorem}{Theorem}
\newtheorem{lemma}{Lemma}
\newcommand{\E}{\mathbb E}
\newcommand{\Prb}{\mathbb P}
\title{Exact blind fiber recovery in random covers of a cycle}
\author{Independent adversarial posterior audit}
\date{November 15, 2026}
\begin{document}
\maketitle

\paragraph{Scope.}
This note independently audits the root's degree-two recovery candidate,
recorded in \path{cycle-cover-recovery-root.txt}.
It proves its exact Bayes success probability for unconditioned random
lifts of a known simple cycle and records the corresponding randomized
minimax formulation. Historical priority is not asserted. The proof is
self-contained; monodromy and the permutation cycle generating function
are standard inputs, both derived below.

\paragraph{Experiment and target.}
Fix integers $k\geq3$ and $N\geq1$, and put $n=kN$.
The known base is the named simple cycle $C_k$ on $\mathbb Z/k\mathbb Z$.
Its $N$-lift initially has $N$ vertices in each fiber and one independent
uniform perfect matching for every base edge. Relabel all $n$ vertices
by an independent uniform permutation, and reveal only the resulting
labeled graph $G$. The target is the \emph{unordered} partition $\Pi$
into the $k$ hidden fibers, not an orientation or a named projection.

An equivalent experiment first chooses a uniform balanced named map
$f:[n]\to\mathbb Z/k\mathbb Z$, with $|f^{-1}(i)|=N$, and then
chooses independent uniform perfect matchings between successive fibers.
The number of such maps is $n!/(N!)^k$. Conditional on $f$, there are
$(N!)^k$ matching choices, all equally likely. The equivalence follows
because a uniform global relabeling gives a uniform balanced named map;
the conditional fiber-wise orderings preserve uniform matching laws.

\begin{theorem}[Exact Bayes and randomized minimax success]
Let $R(G)$ be the number of connected components of $G$. For every graph
in the support of the experiment, the posterior on compatible unordered
fiber partitions is uniform and its support has exactly
\begin{equation}
 M(G)=(2k)^{R(G)-1}.                                      \label{number}
\end{equation}
The optimal conditional Bayes success is $M(G)^{-1}$. The optimal
unconditional success is exactly
\begin{align}
 P_{k,N}
 &=\E(2k)^{1-R(G)}                                      \notag\\
 &=\frac{2k\,\Gamma(N+1/(2k))}
         {\Gamma(1/(2k))\Gamma(N+1)}                    \label{exact}\\
 &\sim\frac{2k}{\Gamma(1/(2k))}\,
                   N^{-1+1/(2k)}\quad(N\to\infty).      \notag
\end{align}
Thus exact blind fiber recovery has success tending to zero even with
unlimited computation. The same expression is the optimal randomized
minimax success when the named balanced hidden projection is an unknown
parameter and success is measured for its unordered fibers.

The unconditioned lift is connected with probability $1/N$. Conditional
on connectedness and with the base $C_k$ supplied, its unordered fiber
partition is unique and can be recovered with probability one. This
conditioning changes the model.
\end{theorem}

\begin{lemma}[Monodromy and component sizes]
There is a uniform permutation $\pi\in S_N$ such that its permutation
cycles of lengths $\ell_1,\ldots,\ell_R$ correspond to the components
of $G$, which are simple cycles of lengths $k\ell_1,\ldots,k\ell_R$.
\end{lemma}
\begin{proof}
Orient the base cycle consistently. Name vertices within the fibers so
that matchings on the first $k-1$ successive edges become identity
matchings: transport names along that spanning path. The last matching
then becomes a permutation $\pi$ on the first fiber. Equivalently $\pi$
is the composition of the original $k$ independent uniform permutations,
with inverses if a different orientation convention is used. It remains
uniform. Following $k$ successive lift edges applies $\pi$, so a cycle
of length $\ell$ in $\pi$ creates one component of length $k\ell$.
Since $k\geq3$ and the base is simple, the resulting graph has no loops
or parallel edges. Every component is a simple cycle. Global vertex
relabeling leaves its component lengths unchanged.
\end{proof}

\begin{lemma}[Cover maps of one component]
A component cycle of length $k\ell$ has exactly $2k$ locally bijective
maps to the named base $C_k$. Every such map puts exactly $\ell$ vertices
in each base fiber.
\end{lemma}
\begin{proof}
Choose the image of one fixed component vertex in $k$ ways, and choose
which of its two neighbors maps in the positive base direction in two
ways. At each later vertex, local bijectivity forces the next edge to
continue in that same direction: reversing would send its two distinct
neighbors to the same base neighbor. Thus the map along a chosen
component orientation has the form $j\mapsto a+j$ or $j\mapsto a-j$
modulo $k$. It closes because the component length is divisible by $k$.
These $2k$ maps are distinct for $k\geq3$ and exhaust all possibilities.
Each base vertex occurs once every $k$ steps, hence exactly $\ell$ times.
This includes $k=3$, where the two base neighbors are still distinct.
\end{proof}

\begin{lemma}[One global automorphism quotient]
For a fixed observed labeled graph with $R$ components, there are
$(2k)^R$ named covering maps to $C_k$, all balanced. Two such maps yield
the same unordered partition if and only if they differ by a single
global automorphism of $C_k$. Each orbit has exactly $2k$ elements.
\end{lemma}
\begin{proof}
Maps on different components may be chosen independently, giving the
product count. Every component contributes the same number of vertices
to every fiber, and its contributions sum to $N$.

Suppose $f,g$ yield the same unordered partition. There is a unique
permutation $\sigma$ of the $k$ base vertices with $g=\sigma\circ f$,
because both maps are surjective and have the same fibers. Every base
edge is the image of a lift edge under $f$. Since $g$ is also a covering
map, $\sigma$ sends each base edge to a base edge. It is therefore a
graph automorphism of the finite cycle. Conversely composing any cover
map with a base automorphism preserves its unordered partition.

The automorphism group of $C_k$ consists of the $k$ rotations and $k$
reflections, so has order $2k$; for $k=3$ this is also $S_3$. Its action
on covering maps is free, since a surjective $f$ can be fixed by
$\sigma\circ f=f$ only when $\sigma$ is the identity. This proves
\eqref{number}.

Repeated component lengths do not change this argument. Observed vertices
are labeled, and exchanging components is not an additional quotient on
partitions of that fixed vertex set. Likewise the independent phase and
orientation choices on components are genuine different partitions
unless absorbed by the \emph{same global} base automorphism.
\end{proof}

\begin{lemma}[Constant likelihood and uniform posterior]
Every compatible pair $(f,G)$ of a balanced named map and an observed
labeled graph has joint probability exactly $1/n!$. Every compatible
unordered partition therefore has joint probability $2k/n!$ with $G$.
\end{lemma}
\begin{proof}
The prior probability of $f$ is $(N!)^k/n!$. Given $f$ and a compatible
$G$, every base-edge matching is determined uniquely by the edges of
$G$ between its two corresponding named fibers. Thus the conditional
likelihood is $(N!)^{-k}$. Their product is $1/n!$.

By the preceding lemma, an unordered partition compatible with $G$ has
exactly $2k$ compatible named projections. Sum their equal joint
probabilities to obtain $2k/n!$, independent of the partition. Dividing
by the marginal probability of $G$ proves the uniform posterior.

In particular the marginal probability of an observed labeled graph is
$(2k)^R/n!$. No division by graph automorphisms appears: $G$ is observed
on a fixed labeled vertex set. The prior on unordered balanced partitions
is uniform as well, but for $k\geq4$ its likelihood mixes different
cyclic orders of its blocks; the named-map calculation handles that
mixture correctly. It is unsafe to assume every labeling of those blocks
is compatible with a fixed named $C_k$.
\end{proof}

\begin{proof}[Proof of the success formula]
No estimator has conditional Bayes success greater than the largest
posterior mass, which is $(2k)^{1-R}$. Choosing any compatible partition
attains that mass. It remains to average over the component count.

Let $R_N$ count the cycles of a uniform permutation of $N$ elements and
put $F_N(t)=\sum_{\pi\in S_N}t^{R_N(\pi)}$. Insert element $N$ into a
permutation of $N-1$ elements. It can be a new singleton cycle, increasing
the cycle count, or be inserted after one of the $N-1$ existing elements,
leaving the count unchanged. Hence
\[
 F_N(t)=(t+N-1)F_{N-1}(t),\qquad F_0(t)=1,
\]
and therefore
\[
 \E t^{R_N}=\frac{t(t+1)\cdots(t+N-1)}{N!}
 =\frac{\Gamma(N+t)}{\Gamma(t)\Gamma(N+1)}\quad(t>0).
\]
Take $t=1/(2k)$ and multiply by $2k$ to obtain \eqref{exact}.
The asymptotic follows from the fixed-parameter gamma ratio
$\Gamma(N+t)/\Gamma(N+1)\sim N^{t-1}$, itself an immediate consequence
of Stirling's formula. Finally $R_N=1$ exactly for the $(N-1)!$ cyclic
permutations, so connectedness has probability $1/N$. At $R=1$,
\eqref{number} is one.
\end{proof}

\paragraph{Fixed parameters and randomized minimax interpretation.}
Let the unknown parameter be any balanced named projection $f$, and let
the observation be drawn by its independent uniform edge matchings.
The preceding Bayes result gives an upper bound on the optimal worst-case
success over $f$, since a worst-case success cannot exceed the average
under its uniform prior.

Conversely, consider an estimator which samples uniformly from all
compatible unordered partitions. For every graph $G$ and every compatible
true $f$, its conditional success over the estimator's own randomization
is exactly $1/M(G)$. For every fixed $f$, the matching monodromy is uniform
on $S_N$, so averaging gives the same value $P_{k,N}$ for every parameter.
It attains the upper bound and proves the randomized minimax assertion.
Sampling such a partition is elementary: orient and start each component
canonically, fix one component's phase and direction, and choose each
remaining component's phase and direction independently uniformly among
its $2k$ choices. Freeness of the global automorphism action proves that
this samples each unordered partition exactly once and uniformly.

The formula is \emph{not} an upper bound on the performance of every
estimator at one specially chosen fixed partition: an estimator may
hardcode that partition and succeed there. The global relabeling prior,
or a worst-case/invariance formulation, is essential to the impossibility
statement. No restriction on computation is used.

\paragraph{Checks and boundaries.}
For $N=1$ the formula gives one. For $N=2$ it gives
$(1+1/(2k))/2$, hence $7/12$ for $k=3$ and $9/16$ for $k=4$.
For $k=3,N=3$ it gives $91/216$. The base cases $k=1,2$ are excluded:
they are not simple cycles with two distinct base neighbors, so the map
count and covering model above do not apply. Connected conditioning gives
perfect recovery for the supplied $C_k$, whereas the unconditioned
success tends to zero. If the base size were unknown, a connected cycle
could cover other cycle bases whose orders divide its length; that is a
different identification question and no uniqueness over those bases
is asserted here.
The result concerns fiber recovery, not recovery of hidden names, a deck
group, or a unique coordinate system within a fiber.

Together with a separately proved recovery theorem for fixed connected
simple regular bases of degree at least three, this supplies a sharp
degree-two boundary for this particular blind, unconditioned random-lift
fiber-recovery problem. It supplies no claim about growing base size,
degree-two models conditioned connected, noisy observations, or priority.

\end{document}
